% !TeX program = xelatex
% =====================================================================
%  示例简历 01 · 后端开发 / 分布式系统 · 本科校招
%  写法说明：只用少量数字，更多写"做了什么、为什么这么做、上线后怎样"。
%  【虚构声明】人物用通用占位名（张三），学校与公司一律以“示例”开头，均为虚构。
%  奖项、指标同样是编造的，不指向任何真实个人、学校或公司。
%  编译：xelatex 01-backend.tex   （或在本仓库 make examples）
% =====================================================================
\documentclass[11pt,a4paper]{article}

\usepackage[UTF8,fontset=none]{ctex}
\usepackage{fontspec}
\usepackage[left=1.15cm,right=1.15cm,top=0.85cm,bottom=0.8cm]{geometry}
\usepackage{enumitem}
\usepackage{tabularx}
\usepackage{titlesec}
\usepackage{xcolor}
\usepackage{graphicx}
\usepackage{microtype}
\usepackage{xurl}
\usepackage{hyperref}

\defaultfontfeatures{
    Extension=.otf,
    UprightFont=*-regular,
    BoldFont=*-bold,
    ItalicFont=*-italic,
    BoldItalicFont=*-bolditalic,
    Numbers=Lining}
\setmainfont{texgyrepagella}
\setsansfont{texgyrepagella}
\defaultfontfeatures{}   % 清空默认项，下面的中文字体各自显式指定
% 正文中文用楷体；楷体无粗体，加粗与标题用黑体，混排更清晰
\IfFontExistsTF{KaiTi}{%
    \IfFontExistsTF{SimHei}{%
        \setCJKmainfont{KaiTi}[BoldFont=SimHei]%
    }{%
        \setCJKmainfont{KaiTi}[AutoFakeBold=2.5]%
    }%
}{%
    \IfFontExistsTF{FandolKai-Regular.otf}{%
        \setCJKmainfont{FandolKai-Regular.otf}[BoldFont=FandolHei-Regular.otf]%
        \setCJKsansfont{FandolHei-Regular.otf}%
    }{%
        \setCJKmainfont{Noto Serif CJK SC}[BoldFont=Noto Sans CJK SC]%
        \setCJKsansfont{Noto Sans CJK SC}%
    }%
}

\definecolor{theme}{HTML}{1A2B3C}       % 主色：沉稳炭青
\definecolor{textmain}{HTML}{222222}    % 正文
\definecolor{textmuted}{HTML}{5A5A5A}   % 次级信息
\definecolor{linecolor}{HTML}{D0D7DE}   % 分割线

\color{textmain}
\pagestyle{empty}
\setlength{\parindent}{0pt}
\setlength{\parskip}{0pt}
\linespread{1.15}
\emergencystretch=2em
\tolerance=3000

% 带下划线的链接（URL 短，不会跨行）
\newcommand{\ulink}[2]{\href{#1}{\color{theme}\underline{#2}}}
\newcommand{\skillitem}[2]{\item \textbf{#1}：\,#2}

% 经历抬头：标题靠左、身份/方向居中、时间靠右（等间距分布，标题过长只会报 Overfull，不会压字）
\newcommand{\entry}[3]{%
    \par\addvspace{2.9pt}%
    \noindent
    \hbox to \linewidth{%
        {\bfseries\fontsize{11.5}{14}\selectfont\color{theme}#1}%
        \hfil{\small\color{textmuted}#2}%
        \hfil{\small\color{textmuted}#3}%
    }\par\nobreak\vspace{1pt}%
}
% 经历下方的一行说明（实习业务方向 / 项目简介）
\newcommand{\subline}[1]{{\small\color{textmuted}#1}\par\nobreak\vspace{1.5pt}}

\titleformat{\section}
    {\large\bfseries\color{theme}}{}{0pt}{}
    [{\vspace{1.5pt}\color{linecolor}\hrule height 0.7pt}]
\titlespacing*{\section}{0pt}{8.4pt}{3pt}

\setlist[itemize]{
    leftmargin=1.15em,
    labelsep=0.45em,
    itemsep=2.9pt,
    parsep=0pt,
    topsep=1.8pt,
    label={\small\color{theme}\textbullet}
}

\hypersetup{
    colorlinks=true,
    urlcolor=theme,
    linkcolor=theme,
    pdfauthor={你的姓名},
    pdftitle={你的姓名 · 个人简历}
}

\begin{document}
\hypersetup{pdfauthor={张三}, pdftitle={张三 · 个人简历}}

\begin{minipage}[c]{0.83\linewidth}
    {\bfseries\fontsize{25}{28}\selectfont\color{theme}张三}%
    \hspace{0.85em}{\fontsize{13.5}{16}\selectfont\color{theme}后端开发 / 分布式系统}\par
    \vspace{4pt}
    {\small\color{textmuted}男 · 22 岁 · 2027 届本科 · 杭州 / 上海}\par
    \vspace{3pt}
    {\small
        138-0000-0001 \quad
        zhangsan@example.com \quad
        \ulink{https://github.com/example}{github.com/example}
    }
\end{minipage}%
\hfill
\begin{minipage}[c]{0.13\linewidth}
    \raggedleft
    {\color{linecolor}\setlength{\fboxsep}{0pt}\setlength{\fboxrule}{0.5pt}%
    \fbox{\parbox[c][2.4cm][c]{1.85cm}{\centering\small\color{textmuted}照片\\（可选）}}}
\end{minipage}

\section{教育背景与荣誉}
\entry{示例理工大学}{计算机科学与技术 · 本科 · 专业排名前 \textbf{8\%}}{2023.09 — 2027.06}
\subline{主修课程：数据结构与算法、操作系统、计算机网络、数据库系统、分布式系统（GPA 3.8 / 4.0）。}
{\small
\begin{tabularx}{\linewidth}{@{}X@{\hspace{1.2em}}X@{}}
    % 荣誉栏不必凑满：这里只放 4 条，两行刚好；不相关的删掉整行即可
    中国大学生计算机设计大赛 \textbf{国家级一等奖} \hfill 2025.08 &
    校级一等奖学金 \hfill 2024.10 \\[3.8pt]
    中国高校计算机大赛 \textbf{省级一等奖} \hfill 2025.04 &
    CET-6 \textbf{560 分} \hfill 2024.12
\end{tabularx}}

\section{实习经历}
\entry{示例科技（杭州）}{后端开发实习生}{2025.06 — 2025.09}
\subline{订单中台结算链路：接口开发、数据一致性与日常值班。}
\begin{itemize}
    \item 负责结算链路的接口开发与维护，把单体接口按读写分离拆成三个服务，热点订单放进 Redis；改造后高峰期每周十几条的接口超时告警基本消失。
    \item 全量对账原本要跑三个多小时，改成按商户分片 + 断点续跑后，凌晨批处理不再和早高峰抢数据库。
    \item 跨库结算的一致性用\textbf{本地消息表 + 定时补偿}保证，配合对账差异告警，上线后没再出现重复扣款的客诉。
    \item 和测试、运维一起梳理结算链路的监控项并接入值班看板，出问题能直接定位到商户和批次。
\end{itemize}

\entry{示例数据}{服务端研发实习生}{2026.01 — 2026.05}
\subline{实时指标平台：数据接入、查询服务与容量保护。}
\begin{itemize}
    \item 参与实时指标平台的数据接入，把 Kafka 消费改成批量落库并做削峰，下游看板不再出现数据回退。
    \item 用 \textbf{ClickHouse 预聚合 + 结果缓存}改写查询服务，原本要等几秒的报表回到几百毫秒，运营同学不用提前一天跑数。
    \item 给上报接口加了\textbf{令牌桶限流}与优先级队列，突发流量下优先保核心查询，大促期间没有做过人工扩容。
\end{itemize}

\section{项目经历}
\entry{Lattice — 轻量 RPC 框架}{独立开发 · C++17 · \ulink{https://github.com/example/lattice}{github.com/example/lattice}}{2025.10 — 2026.02}
\subline{自己写的 RPC 框架：多 Reactor 网络模型 + 自定义二进制协议，带服务治理能力。}
\begin{itemize}
    \item 网络层用 \textbf{epoll 多 Reactor} 配合无锁队列；为了调通边缘触发和线程模型，压测脚本改了四五版才稳定。
    \item 服务治理做了注册发现、一致性哈希负载均衡与超时重试；节点被摘除时客户端会自动切换，不会把请求继续打到坏节点上。
    \item 用 GTest 补齐协议编解码、连接管理与重试逻辑的用例，配 GitHub Actions 做构建和冒烟压测。
\end{itemize}

\entry{KeyLite — KV 存储引擎}{独立开发 · Rust · LSM-Tree}{2026.03 — 2026.07}
\begin{itemize}
    \item 按 \textbf{LSM-Tree} 的分层思路实现写入路径，用 WAL 保证进程崩溃后能恢复到最后一条已确认的写入。
    \item 内存索引用跳表加布隆过滤器，点查不用扫盘；另外写了故障注入脚本验证崩溃恢复流程。
\end{itemize}

\section{专业技能}
\begin{itemize}
    \skillitem{编程语言}{C++（C++17 / 20）、Go、Python、SQL}
    \skillitem{后端与中间件}{MySQL、Redis、Kafka、ClickHouse、gRPC、分布式锁与限流}
    \skillitem{系统与工程}{Linux、多线程与网络编程、Docker、Git、CMake、性能分析}
    \skillitem{协作与其他}{代码评审、技术方案文档、跨端联调与值班；英文技术文档阅读}
\end{itemize}

\end{document}
